\paragraph{Proposition15 (mixing implies spectral).} There exists a universal constant $c_{\mathrm{spec}}>0$ for which the following holds. Let $H=(V=[n],E)$ be any $d$-regular hypergraph, and $\lambda\leq1/2$. If $H$ is $\lambda$-mixing and symmetric, then $H$ is a $\lambda'$-spectral for $\lambda'=c_{\mathrm{spec}}\cdot\lambda\log(1/\lambda)$.
The proof of Proposition15 is a simple application of a similar result for graphs.
\paragraph{Lemma16 (Lemma3.3 of[7]).} There exists a universal constant $c_{\mathrm{BL}}$ for which the following hods. For any $n\times n$ real symmetric matrix $A$ and $\lambda\leq1/2$, suppose each row has an $\ell_1$ norm of at most1 and for any two vectors $u,v\in\{0,1\}^n$,\footnote{Note that Bilu and Linial's Lemma statement only has the weaker requirement that this hold for orthogonal $u$ and $v$. As a result, they have an additional condition that the diagonal entries of $A$ not be too large, which is not needed for our version of the Lemma.}
\[|u^{\dagger}Av|\leq\lambda\|u\|\|v\|.\tag{4}\]
Then, the spectral radius of $A$ is at most $c_{\mathrm{BL}}\cdot\lambda\log(1/\lambda)$.
If $G$ is the (normalized) transition matrix of a $d$-regular graph, and $A=G-\frac1nJ$ for $J$ being the all-ones matrix, then Lemma16 shows a converse to the expander mixing lemma: Any graph with good mixing is also a good spectral expander. We show that their result can be lifted to3-uniform hypergraphs.
\begin{proof}[Proof of Proposition15]
Fix any $y\in\mathbb R^n$ and let $A^{(y)}\in\mathbb R^{n\times n}$ be defined as
\[A^{(y)}_{i,k}\coloneqq\frac1{2d}\left(\sum_{j\in[n]}1[(i,j,k)\in E]\cdot y_j\right)-\frac1{2n^2}\sum_{j\in[n]}y_j.\]
Then, for any $x,z\in\mathbb R^n$,
\[2(z^{\dagger}A^{(y)}x)=\frac1d\sum_{(i,j,k)\in E}x_iy_jz_k-\frac1{n^2}\sum_{i\in V}x_i\sum_{i\in V}y_i\sum_{i\in V}z_j.\]
Therefore, in order to prove that $H$ is an $\lambda'$-spectral expander, it is sufficient to show that for all $y\in\mathbb R^n$, the operator norm of $A^{(y)}$ is at most $\|y\|_{\infty}\cdot\lambda'/2$. We observe that:
\begin{enumerate}
\item For any fixed $x,z\in\mathbb R^n$, the quantity $z^{\dagger}A^{(y)}x$ is a linear function of $y$. Thus, we can consider $y$-s with $\|y\|_{\infty}=1$ without loss of generality. Furthermore, the $y$ maximizing $z^{\dagger}A^{(y)}x$ among those with $\|y\|_{\infty}=1$ will be in $\{\pm1\}^n$. Therefore we assume $y\in\{\pm1\}^n$ also without loss of generality.
\item Since $H$ is symmetric, the operator norm and spectral radius of $A^{(y)}$ are equal, so we instead bound the spectral radius.
\end{enumerate}
We will apply Lemma16 to each $A^{(y)}$. Note that:
\begin{itemize}
\item $A^{(y)}$ is symmetric, which follows immediately from the fact that $H$ is symmetric.
\item The $\ell_1$ norm of each row of $A^{(y)}$ is bounded by1:
\begin{multline*}
\sum_{k\in[n]}\left|\frac1{2d}\left(\sum_{j\in[n]}1[(i,j,k)\in E]\cdot y_j\right)+\frac1{2n^2}\sum_{j\in[n]}y_j\right|\\
\leq\frac1{2d}\cdot d+\frac1{2n^2}\cdot n\leq1.
\end{multline*}
\item Finally, fix some $u,v\in\{0,1\}^n$, and define the sets
\[S_1=\{i\in V:u_i=1\},\quad S_3=\{i\in V:v_i=1\},\quad S_2^+=\{i\in V:y_i=1\},\quad S_2^-=\{i\in V:y_i=-1\}.\]
Then,
\begin{align*}
|u^{\dagger}A^{(y)}v|&=\frac12\left|\frac1d(E(S_1,S_2^+,S_3)-E(S_1,S_2^-,S_3))\right.\\
&\hspace{3em}\left.-\frac{|S_1|\cdot|S_3|\cdot(|S_2^+|-|S_2^-|)}{n^2}\right|\\
&\leq\frac12\left|\frac1dE(S_1,S_2^+,S_3)-\frac{|S_1|\cdot|S_3|\cdot|S_2^+|}{n^2}\right|\\
&\quad+\frac12\left|\frac1dE(S_1,S_2^-,S_3)-\frac{|S_1|\cdot|S_3|\cdot|S_2^-|}{n^2}\right|\\
&\leq\lambda\sqrt{|S_1|\cdot|S_3|},
\end{align*}
where the last inequality follows from fact that $H$ is $\lambda$-mixing (applied to both expressions).
\end{itemize}
Thus, $|u^{\dagger}A^{(y)}v|\leq\lambda\|u\|\|v\|$ and we can apply Lemma16 to obtain $\|A^{(y)}\|_{\mathrm{op}}=c_{\mathrm{BL}}\cdot\lambda\log(1/\lambda)$, implying that $H$ is $(\lambda'=2c_{\mathrm{BL}}\cdot\lambda\log(1/\lambda))$-spectral.
\end{proof}
